% ===== Chapter 4 opening: what makes the human brain special (~3 min) =====
\partopener{Գլուխ 4}{Խելք}{Կարծում ես՝ խելացի՞ ես։ Խելքի շփոթեցնող գիտությունը}{smug}{peach}

% The chapter's first paragraph: philosophy, vehicles ... and 300 kinds of pasta
\begin{frame}{Ինչով է մարդու ուղեղը յուրահատուկ}
  \begin{columns}[c]
    \begin{column}{0.55\textwidth}
      Շնչել, ուտել, քնել ու վախենալ կարողանում են շատ կենդանիներ։ Բայց ոչ մի այլ արարած չի ստեղծել իր...\par\medskip
      \begin{bookquote}
        ...փիլիսոփայությունը, տրանսպորտը, հագուստը, էներգիայի աղբյուրները, կրոնը կամ \textbf{մակարոնի թեկուզ մեկ տեսակ}, էլ չասած՝ երեք հարյուրից ավելի։
      \end{bookquote}
    \end{column}
    \begin{column}{0.4\textwidth}
      \centering
      \bigemoji[1.1cm]{thinking}\ \bigemoji[1.1cm]{car}\ \bigemoji[1.1cm]{jeans}\par\smallskip
      \bigemoji[1.1cm]{lightning}\ \bigemoji[1.1cm]{pasta}\par\medskip
      \onslide<2->{%
      \begin{card}[colback=peach!40!cardwhite]
        Գիրքը հիմնականում ուղեղի թերությունների մասին է, բայց ինչ-որ բան նա \textbf{հաստատ ճիշտ է անում}։
      \end{card}}
    \end{column}
  \end{columns}
\end{frame}

% The famous aphorism and Burnett's circular-reasoning joke
\begin{frame}{Ուղեղը, որ փորձում է հասկանալ ինքն իրեն}
  \begin{card}[colback=lavender!35!cardwhite, fontupper=\mediumfont]
    \centering «Եթե մարդու ուղեղն այնքան պարզ լիներ, որ կարողանայինք այն հասկանալ, մենք այնքան պարզ կլինեինք, որ չէինք կարողանա»\par
    {\raggedleft\scriptsize\color{plum!70!cream} հայտնի աֆորիզմ\par}
  \end{card}
  \vspace{2mm}
  \begin{columns}[c]
    \begin{column}{0.62\textwidth}
      \onslide<2->{%
      Գիտությունը կա մեր խելքի շնորհիվ, իսկ հիմա գիտությամբ ուսումնասիրում ենք հենց խելքը։\par\smallskip}
      \onslide<3->{%
      \begin{bookquote}
        Սա կա՛մ շատ արդյունավետ է, կա՛մ շրջանաձև տրամաբանություն։ Ես բավականաչափ խելացի չեմ, որ որոշեմ։
      \end{bookquote}}
    \end{column}
    \begin{column}{0.32\textwidth}
      \centering
      \onslide<3->{\mascot[2.2cm]{confused}}
    \end{column}
  \end{columns}
\end{frame}

% Chat: one word for "intelligence", then Burnett's description of it
\begin{frame}{Ի՞նչ է խելքը}
  \begin{columns}[c]
    \begin{column}{0.47\textwidth}
      \begin{chatbox}
        Գրիր \textbf{մեկ բառով}. ի՞նչ է խելքը։\par\smallskip
        {\footnotesize Օրինակ՝ «արագություն», «գիտելիք», «հնարամտություն»...}
      \end{chatbox}
    \end{column}
    \begin{column}{0.47\textwidth}
      \onslide<2->{%
      \begin{card}[colback=peach!40!cardwhite]
        {\blackfont Շփոթեցնող, խառը, հաճախ հակասական ու դժվար ըմբռնելի}\par\smallskip
        Բըրնեթի խոսքով՝ սա \textbf{հենց խելքի} լավագույն նկարագրությունն է։
      \end{card}}
    \end{column}
  \end{columns}
  \vspace{4mm}
  \onslide<3->{\centering Այսօր՝ ինչպես ենք օգտագործում խելքը, ինչու է այն դժվար չափել, ու ինչ տարօրինակ հատկություններ ունի։\par}
\end{frame}
